\section{Program-Verify Loops and Write Latency}
\label{sec:bg-pvloops}

Because OxRAM switching is stochastic, a single programming pulse cannot
reliably bring a cell to a target resistance state. Commercial
architectures therefore employ a program-verify (PV) loop: a sequence of
programming pulses interleaved with resistance measurements that
continues until the cell resistance falls within a target
window~\cite{kawahara2013filament, milo2021accurate}. Each iteration of
the loop applies a pulse and checks the result; if the cell has not
reached the target state, another pulse is applied. The loop terminates
when the cell resistance is confirmed to be within specification.
Single-level cells (SLCs) define two resistance windows corresponding to
logical 0 and 1, while multi-level cells (MLCs) require finer resistance
windows and more sophisticated verification schemes~\cite{hung2021mlc}.

The completion time of a PV loop, and therefore the write latency $t_w$
of the device, is non-deterministic. Two sources of variability
contribute to $t_w$. The stochastic component arises from cycle-to-cycle
variability in filament switching dynamics~\cite{6582090, 6724732}: the
number of pulses required to grow or rupture a filament varies even
under identical initial conditions. The deterministic component arises
because the number of PV iterations required to reach the target
resistance scales with the distance between the initial cell resistance
and the target threshold~\cite{yu2012switching, kawahara2013filament}. A
transition originating far from the target requires more programming
pulses than one originating near it. Together, these contributions make
write latency a function of both the stochastic device physics and the
logical state being written.

The data-dependence of write latency has direct implications for timing
side-channel analysis. Writes that require only SET transitions, only
RESET transitions, or no transitions at all produce different latency
distributions~\cite{yu2012switching}. In a crossbar architecture, writes
that require both SET and RESET transitions within the same word incur
additional serialization overhead, as described in
Section~\ref{sec:bg-crossbar}. The presence of an on-die ECC further
modifies the effective transition composition of every write, since
parity bits computed over user data may require transitions that differ
in polarity from the data transitions. The net result is that the write
latency of a commercial ReRAM device encodes information about the
codeword being written, not merely the user-visible data word. This
relationship is the basis of the BREW-RC technique presented in
Chapter~\ref{chap:brewrc}.